\chapter{埃尔米特型}\label{ch:b2:hermitian}

复向量空间有其自身的内积几何，只是多了一处不同：
一个变量上线性，另一个变量上\emph{共轭线性}。
接受这一不同所带来的回报，是比实数情形更简洁的谱理论：
\omterm{def:b2:hermitian:adjoint}{埃尔米特自同态}的\omterm{def:b2:reduction:eigen}{特征值}均为实数，
\omterm{def:b2:hermitian:adjoint}{酉自同态}的\omterm{def:b2:reduction:eigen}{特征值}模长均为 1，两者都可在标准正交基下对角化。
本章较短，
将\cref{ch:b2:quadratic}中的欧几里得理论平移到 $\C$ 上。

\section{埃尔米特内积}

\begin{definition}\label{def:b2:hermitian:def}
\emph{埃尔米特内积}\index{埃尔米特内积}是复向量空间
$E$ 上的映射 $\langle\cdot,\cdot\rangle
\colon E \times E \to \C$，对第二个变元线性，
满足共轭对称性（$\langle y, x\rangle = \conj{\langle x,
y\rangle}$——因而对第一个变元共轭线性），并且
正定（$\langle x, x\rangle > 0$，当 $x \neq 0$ 时）。
在 $\C^n$ 上的标准例子为：
\[
\langle x, y \rangle = \sum_{i=1}^{n} \conj{x_i}\, y_i ;
\]
对\omterm{def:b2:metric:continuity}{连续}函数，$\langle f, g\rangle = \int_a^b \conj
f\,g$。\omterm{def:b2:nvs:norm}{范数} $\norm x = \sqrt{\langle x,x\rangle}$；有限维复
向量空间若如此赋予内积，则称为\emph{埃尔米特
空间}\index{埃尔米特空间}。
\end{definition}

\begin{example}[初步计算]\label{ex:b2:hermitian:firstcomp}
在 $\C^2$ 中取 $x = (1+\iu,\ 2-\iu)$ 与 $y = (\iu,\ 1)$。
于是有
\[
\norm x^2 = \abs{1+\iu}^2 + \abs{2-\iu}^2 = 2 + 5 = 7,
\qquad
\norm y^2 = 1 + 1 = 2,
\]
并对第一个变元取共轭，
\[
\langle x, y\rangle
= \conj{(1+\iu)}\,\iu + \conj{(2-\iu)}\cdot1
= (1-\iu)\iu + (2+\iu)
= (\iu + 1) + (2 + \iu) = 3 + 2\iu .
\]
柯西–施瓦茨不等式成立：$\abs{\langle x, y\rangle}^2 = 9 +
4 = 13 \leq 14 = \norm x^2\norm y^2$——接近取等，
因为 $x$ 接近 $y$ 的某个倍数。另外注意
$\langle y, x\rangle = \conj{3 + 2\iu} = 3 - 2\iu$:
这正是共轭对称性的体现，也说明了为何 $\langle x,
x\rangle$ 总是实数。
\end{example}

\begin{theorem}[柯西–施瓦茨不等式（复数情形）]\label{thm:b2:hermitian:cs}
$\abs{\langle x, y\rangle} \leq \norm x \norm y$，等号成立当且仅当
$x, y$ 线性相关；$\norm\cdot$ 为\omterm{def:b2:nvs:norm}{范数}。此外
存在标准正交基（施密特正交化原样适用），且有
\[
x = \sum_i \langle e_i, x\rangle\, e_i,
\qquad
\norm x^2 = \sum_i \abs{\langle e_i, x\rangle}^2 .
\]
\end{theorem}

\begin{proof}
对 $y \neq 0$ 及 $t \in \C$，有 $0 \leq \norm{x - ty}^2 = \norm x^2
- 2\Re\bigl(\conj t\langle y, x\rangle\bigr) + \abs t^2\norm y^2$.
取 $t = \frac{\langle y, x\rangle}{\norm y^2}$：
\[
0 \leq \norm x^2 - \frac{\abs{\langle y, x\rangle}^2}{\norm y^2},
\]
此即所证不等式；等号成立时必有 $x = ty$。三角不等式
完整写出如下：
\[
\norm{x + y}^2 = \norm x^2 + 2\,\Re\langle x, y\rangle
+ \norm y^2
\leq \norm x^2 + 2\,\abs{\langle x, y\rangle} + \norm y^2
\leq \bigl(\norm x + \norm y\bigr)^2 ,
\]
利用 $\Re z \leq \abs z$，再用柯西–施瓦茨不等式；齐次性
与正定性显然成立，故 $\norm\cdot$ 为\omterm{def:b2:nvs:norm}{范数}。
施密特正交化：与实情形相同，共轭的
位置按定义确定（注意约定：本书的内积对\emph{第一}个变元
共轭线性，故坐标为
$\langle e_i, x\rangle$；下一例将
完整地运行一次该算法）。
\end{proof}

\begin{example}[复数情形的施密特正交化，完整演算]\label{ex:b2:hermitian:gramschmidt}
将基 $v_1 = (1, \iu)$，$v_2 = (0, 1)$ 标准正交化，
它们是 $\C^2$ 的基。第一个向量：$\norm{v_1}^2 = \abs1^2 + \abs\iu^2 = 2$，
故 $e_1 = \frac{1}{\sqrt2}(1, \iu)$。对 $v_2$ 作投影——注意
在第一个变元上取共轭：
\[
\langle e_1, v_2\rangle
= \frac{1}{\sqrt2}\bigl(\conj{1}\cdot0 +
\conj{\iu}\cdot1\bigr)
= \frac{-\iu}{\sqrt2} ,
\qquad
v_2 - \langle e_1, v_2\rangle e_1
= (0,1) + \frac{\iu}{2}\,(1, \iu)
= \Bigl(\frac\iu2,\ \frac12\Bigr) .
\]
其\omterm{def:b2:nvs:norm}{范数}为 $\sqrt{\frac14 + \frac14} = \frac{1}{\sqrt2}$：
$e_2 = \frac{1}{\sqrt2}(\iu, 1)$。验证：$\langle e_1,
e_2\rangle = \frac12(\conj1\cdot\iu + \conj\iu\cdot1) =
\frac12(\iu - \iu) = 0$。$v_2$ 在新基下的坐标为：
$v_2 = \langle e_1, v_2\rangle e_1 + \langle e_2, v_2\rangle
e_2$，其中 $\langle e_2, v_2\rangle = \frac{1}{\sqrt2}$——
注意次序：$\langle v_2, e_1\rangle$ 给出的是
\emph{共轭}系数。结论：该算法与欧几里得情形完全相同，
唯一的陷阱在于
共轭落在何处，而计算 $\norm{v_2 -
\text{投影}}^2 > 0$ 悄然用到了正定性——这条
使整个几何成立的公理。
\end{example}

\section{伴随、埃尔米特与酉自同态}

\begin{definition}\label{def:b2:hermitian:adjoint}
$u^*$ 是 $u \in \mathcal{L}(E)$ 的\emph{伴随}，由下式定义：
$\langle u^*(x), y\rangle = \langle x, u(y)\rangle$；在
标准正交基下，$\operatorname{Mat}(u^*) = \conj{A}^{\mathsf T} =:
A^{\dagger}$（共轭\omterm{def:b2:linalg:transpose}{转置}）——事实上，若 $B =
(b_{ij})$ 为 $u^*$ 在标准正交基
$(e_i)$ 下的矩阵，则 $b_{ij} = \langle e_i, u^*(e_j)\rangle$，由
定义式可得
\[
\conj{b_{ij}}
= \langle u^*(e_j), e_i\rangle
= \langle e_j, u(e_i)\rangle = a_{ji} ,
\qquad\text{于是}\qquad
B = \conj{A}^{\mathsf T} .
\] 称 $u$ 为
\emph{埃尔米特自同态}\index{埃尔米特自同态}，当 $u^* = u$
（$A^\dagger = A$），为\emph{酉自同态}\index{酉群}，当 $u^*u =
\mathrm{id}$（$A^\dagger A = I$：群 $U(n)$），为\emph{正规}的，
当 $u^*u = uu^*$ 时。
\end{definition}

\begin{proposition}\label{prop:b2:hermitian:eigenvalues}
\omterm{def:b2:hermitian:adjoint}{埃尔米特自同态}的\omterm{def:b2:reduction:eigen}{特征值}都是\emph{实数}；
\omterm{def:b2:hermitian:adjoint}{酉自同态}的\omterm{def:b2:reduction:eigen}{特征值}都有\emph{模为 $1$}；
两种情形下，不同\omterm{def:b2:reduction:eigen}{特征值}所对应的\omterm{def:b2:reduction:eigen}{特征子空间}彼此正交。
\end{proposition}

\begin{proof}
埃尔米特情形，设 $u(x) = \lambda x$，$x \neq 0$：
\[
\lambda \norm x^2 = \langle x, u(x)\rangle = \langle u(x), x\rangle
= \conj{\langle x, u(x)\rangle} = \conj\lambda\,\norm x^2 ,
\]
于是 $\lambda \in \R$。酉情形：$\norm{u(x)} = \norm x$（由 $u^*u =
\mathrm{id}$），故 $\abs\lambda\norm x = \norm x$。正交性
（埃尔米特情形）：$\lambda\langle x, y\rangle = \langle u(x),
y\rangle = \langle x, u(y)\rangle = \mu\langle x, y\rangle$，对
实\omterm{def:b2:reduction:eigen}{特征值} $\lambda \neq \mu$ 的\omterm{def:b2:reduction:eigen}{特征向量}成立。酉情形，完整地
写出：设 $u(x) = \lambda x$，$u(y) = \mu y$，且 $\lambda \neq
\mu$（两者的模均为 1），
\[
\langle x, y\rangle = \langle u(x), u(y)\rangle
= \conj\lambda\mu\,\langle x, y\rangle ,
\]
且 $\conj\lambda\mu = \frac{\mu}{\lambda} \neq 1$，该因子
不等于 $1$，故 $\langle x, y\rangle = 0$。
\end{proof}

\begin{example}[反埃尔米特矩阵的对角化]\label{ex:b2:hermitian:skewexample}
$A = \begin{pmatrix} 0 & -2\\ 2 & 0\end{pmatrix}$ 满足
$A^\dagger = A^{\mathsf T} = -A$：反埃尔米特矩阵（同时也是实
反对称矩阵——在 $\R$ 上它根本没有\omterm{def:b2:reduction:eigen}{特征值}）。
\omterm{def:b2:reduction:charpoly}{特征多项式} $X^2 + 4$：\omterm{def:b2:reduction:eigen}{特征值}为 $\pm2\iu$，
纯虚数，正如\cref{exo:b2:hermitian:9}在
一般情形下所预言的那样。\omterm{def:b2:reduction:eigen}{特征向量}：$(A - 2\iu I)v = 0$ 给出 $v_1 =
\frac{1}{\sqrt2}(1, \iu)$，且 $v_2 = \frac{1}{\sqrt2}(1,
-\iu)$ 对应 $-2\iu$；它们彼此正交：
\[
\langle v_1, v_2\rangle
= \tfrac12\bigl(\conj{1}\cdot1 +
\conj{\iu}\cdot(-\iu)\bigr)
= \tfrac12(1 - 1) = 0 .
\]
于是 $A = U\operatorname{diag}(2\iu, -2\iu)\,U^\dagger$，其中 $U =
(v_1\ v_2)$ 为酉矩阵。最后的观察：$H = -\iu A =
\begin{pmatrix} 0 & 2\iu\\ -2\iu & 0\end{pmatrix}$ 是埃尔米特矩阵，
具有实谱 $\{\pm2\}$，并具有\emph{相同}的
\omterm{def:b2:reduction:eigen}{特征向量}——映射 $u \mapsto \iu u$ 是
\omterm{def:b2:hermitian:adjoint}{埃尔米特自同态}与反\omterm{def:b2:hermitian:adjoint}{埃尔米特自同态}之间的双射
（见\cref{exo:b2:hermitian:9}），从矩阵角度看；在 $\R$ 上
同一个 $A$ 是旋转伸缩变换，根本没有
\omterm{def:b2:reduction:eigen}{特征向量}，只有过渡到 $\C$ 才显露其标准形。
\end{example}

\begin{theorem}[埃尔米特谱定理]\label{thm:b2:hermitian:spectral}
\omterm{def:b2:hermitian:def}{埃尔米特空间}上的每个\omterm{def:b2:hermitian:adjoint}{埃尔米特自同态}都具有一组
由\omterm{def:b2:reduction:eigen}{特征向量}构成的标准正交基（\omterm{def:b2:reduction:eigen}{特征值}均为实数）：$A^\dagger = A$
蕴含 $A = U D U^{\dagger}$，其中 $U \in U(n)$，$D$ 为实
对角矩阵。\index{谱定理!埃尔米特}
\end{theorem}

\begin{proof}
在 $\C$ 上，\omterm{def:b2:reduction:charpoly}{特征多项式}完全分裂，因此有一个\omterm{def:b2:reduction:eigen}{特征向量}
$e_1$ 存在（见\cref{ch:b2:reduction}）——无需紧性论证，
这是 $\C$ 的一个优点。将其单位化。其正交补
$F = e_1^\perp$ 是稳定的：对 $x \perp e_1$，
\[
\langle e_1, u(x)\rangle = \langle u(e_1), x\rangle
= \lambda_1\langle e_1, x\rangle = 0
\]
（$\lambda_1$ 为实数）。限制仍是埃尔米特的；对
维数归纳，再拼接。具体地说，限制 $u|_F$ 是
\omterm{def:b2:hermitian:def}{埃尔米特空间} $F$ 上的一个自同态（维数为 $n - 1$），
它满足 $\langle u|_F(x), y\rangle = \langle x, u|_F(y)\rangle$，
由 $u$ 继承而来；归纳假设给出一组
由\omterm{def:b2:reduction:eigen}{特征向量}构成的标准正交基 $(e_2, \dots, e_n)$（属于 $F$），
于是 $(e_1, e_2, \dots, e_n)$ 在 $E$ 中是标准正交的（$e_1 \perp
F$），并且由 $u$ 的\omterm{def:b2:reduction:eigen}{特征向量}构成。矩阵的表述：
$U$ 的各列即 $e_i$，$U^\dagger U = I$ 表示它们的
标准正交性，$AU = UD$ 汇总了\omterm{def:b2:reduction:eigen}{特征值}
方程，由此得 $A = UDU^\dagger$，其中 $D$ 为实对角
矩阵（见\cref{prop:b2:hermitian:eigenvalues}）。
\end{proof}

\begin{example}\label{ex:b2:hermitian:pauli}
$A = \begin{pmatrix} 0 & -\iu\\ \iu & 0\end{pmatrix}$ 是埃尔米特矩阵
（$A^\dagger = A$）：由 $\chi_A = X^2 - 1$ 得\omterm{def:b2:reduction:eigen}{特征值} $\pm 1$
（实数，正如所预告的），附带标准正交的\omterm{def:b2:reduction:eigen}{特征向量}
$\frac{1}{\sqrt2}(1, \iu)^{\mathsf T}$ 与 $\frac{1}{\sqrt2}(1,
-\iu)^{\mathsf T}$。（物理学家称 $A$ 为泡利矩阵；
埃尔米特谱恒为实数，这正是量子可观测量被
埃尔米特算子刻画的原因。）
\end{example}

\begin{example}[正定埃尔米特矩阵的计算]\label{ex:b2:hermitian:pdrun}
$A = \begin{pmatrix} 2 & 1-\iu\\ 1+\iu & 3\end{pmatrix}$:
是埃尔米特矩阵：对角元为实数，非对角元互为共轭，
\omterm{def:b2:reduction:charpoly}{特征多项式}为
\[
(2-\lambda)(3-\lambda) - \abs{1-\iu}^2
= \lambda^2 - 5\lambda + 4
= (\lambda - 1)(\lambda - 4) :
\]
谱为 $\{1, 4\}$，均为正实数——故 $A$ 为正定。
\omterm{def:b2:reduction:eigen}{特征向量}：对 $\lambda = 4$，方程组 $(A -
4I)v = 0$ 给出 $v_4 = (1 - \iu,\ 2)$（检验第二行：
$(1+\iu)(1-\iu) - 2 = 0$）；对 $\lambda = 1$，$v_1 = (1 - \iu,\
-1)$。正交性（共轭在第一个位置）：
\[
\langle v_4, v_1\rangle
= \conj{(1-\iu)}\,(1-\iu) + \conj{2}\,(-1)
= 2 - 2 = 0 . \checkmark
\]
单位化（$\norm{v_4}^2 = 2 + 4 = 6$，$\norm{v_1}^2 = 2 + 1
= 3$）得酉矩阵 $U = \bigl(\frac{v_4}{\sqrt6}\
\frac{v_1}{\sqrt3}\bigr)$，使 $A =
U\operatorname{diag}(4,1)U^\dagger$。小结：
由\omterm{pb:b2:hermitian:1}{瑞利商}的观点可立即看出——在 $\C^2$ 的单位球面上，
$\langle x, Ax\rangle$ 的值遍及 $\intcc{1}{4}$，
在两个\omterm{def:b2:reduction:eigen}{特征向量}处取到，这是在周末问题中建立的
柯朗–费希尔极小极大定理在 $n = 2$ 时的雏形。
再作一个检验：在非\omterm{def:b2:reduction:eigen}{特征向量}处，该形式的值也是实数。
取 $x = (1, \iu)$，
\[
Ax = \bigl(2 + (1-\iu)\iu,\ (1+\iu) + 3\iu\bigr)
= (3 + \iu,\ 1 + 4\iu),
\]
\[
\langle x, Ax\rangle = \conj{1}\,(3+\iu) +
\conj{\iu}\,(1+4\iu)
= (3 + \iu) + (-\iu)(1 + 4\iu) = 3 + \iu - \iu + 4 = 7 \in
\R ,
\]
正如 \cref{prop:b2:hermitian:eigenvalues} 证明中的机制
（即共轭对称性配合 $A^\dagger = A$）
对每个 $x$ 都保证其为实数。
\end{example}

\begin{example}[酉矩阵的对角化]\label{ex:b2:hermitian:unitary}
$U = \frac{1}{\sqrt2}\begin{pmatrix} 1 & \iu\\ \iu & 1
\end{pmatrix}$ 是酉矩阵（见 \cref{exo:b2:hermitian:2}）。其
\omterm{def:b2:reduction:charpoly}{特征多项式}为 $\bigl(X - \frac{1}{\sqrt2}\bigr)^2
+ \frac12$，根为
\[
\lambda_\pm = \frac{1 \pm \iu}{\sqrt2} = \eu^{\pm\iu\pi/4},
\]
两根的模均为 1，正如 \cref{prop:b2:hermitian:eigenvalues} 所预告，
并有标准正交的\omterm{def:b2:reduction:eigen}{特征向量} $\frac{1}{\sqrt2}(1, \pm1)$。于是 $U =
V\operatorname{diag}(\eu^{\iu\pi/4},
\eu^{-\iu\pi/4})V^\dagger$：在合适的基下，$U$ 是一对
平面旋转，角为 $\pm\frac\pi4$——实旋转矩阵
没有实\omterm{def:b2:reduction:eigen}{特征向量}，但在 $\C$ 上它分裂为两个
模为 1 的纯量。小结：埃尔米特谱位于实直线上，
酉谱位于单位圆上；两者
都是同一正规性的影子，而
\cref{exo:b2:hermitian:6} 中的凯莱变换将一幅图景映为另一幅。
\end{example}

\begin{example}[凯莱变换的计算]\label{ex:b2:hermitian:cayleyrun}
对 $H = \begin{pmatrix} 0 & 1\\
1 & 0\end{pmatrix}$ 运用 \cref{exo:b2:hermitian:6}（埃尔米特，谱为 $\{1, -1\}$，
标准正交的\omterm{def:b2:reduction:eigen}{特征向量}为 $\frac{1}{\sqrt2}(1, \pm1)$）。在特征基下一切化为纯量运算：
变换 $\lambda
\mapsto \frac{\lambda - \iu}{\lambda + \iu}$ 把
\[
1 \longmapsto \frac{1 - \iu}{1 + \iu} = -\iu,
\qquad
-1 \longmapsto \frac{-1 - \iu}{-1 + \iu} = \iu ,
\]
分别映到上式右端的值（分子分母同乘分母的共轭），于是 $U = (H -
\iu I)(H + \iu I)^{-1}$ 为酉矩阵，其在这些相同\omterm{def:b2:reduction:eigen}{特征向量}上的\omterm{def:b2:reduction:eigen}{特征值}为
$\mp\iu$：
\[
U = \frac{1}{2}\begin{pmatrix} 1 & 1\\ 1 & -1\end{pmatrix}
\begin{pmatrix} -\iu & 0\\ 0 & \iu\end{pmatrix}
\begin{pmatrix} 1 & 1\\ 1 & -1\end{pmatrix}
= \begin{pmatrix} 0 & -\iu\\ -\iu & 0\end{pmatrix} .
\]
验证：$U^\dagger U = I$，且 $1 \notin \operatorname{Sp}U =
\{\pm\iu\}$，正如理论所预告的。小结：
实直线映到去掉点 $1$ 的单位圆上——
逐个\omterm{def:b2:reduction:eigen}{特征值}看，凯莱变换即
莫比乌斯映射 $\frac{\lambda-\iu}{\lambda+\iu}$，而矩阵
则只是跟随其谱而变化。
\end{example}

\begin{remark}[常见误区]
\emph{(i) 共轭号落在何处：}本书对
\emph{第一个}位置取共轭，故坐标为 $\langle e_i, x\rangle$，
并有 $\langle\lambda x, y\rangle = \conj\lambda\langle x,
y\rangle$；许多教材则改为对第二个位置取共轭——
比较公式前须先统一约定，否则 $\iu$ 的符号会
悄然出错。\emph{(ii) 复情形下的极化更强：}
在 $\C$ 上，若 $\langle x, u(x)\rangle = 0$ 对\emph{所有} $x$ 成立，
则 $u = 0$（展开 $x + y$ 与 $x + \iu y$，使
$\langle x, u(y)\rangle$ 的实部与虚部皆为零）；
在 $\R$ 上则不成立——旋转 $\frac\pi2$ 满足
$\langle x, u(x)\rangle = 0$ 处处成立。因此在
仅在 $\C$ 上，“$\langle x, u(x)\rangle \in \R$ 对一切 $x$ 成立”
已足以迫使 $u$ 为埃尔米特。\emph{(iii) 实正规不等于
\omterm{def:b2:reduction:diag}{可对角化}：}
\cref{ex:b2:hermitian:skewexample} 中的矩阵是正规的，但没有实\omterm{def:b2:reduction:eigen}{特征值}；
酉对角化是\emph{在}
$\C$ 上的定理，而在 $\R$ 上只能得到分块约化。
\emph{(iv) 检验酉性：}$U^\dagger U = I$ 表示
\emph{列}向量关于\omterm{def:b2:hermitian:def}{埃尔米特内积}标准正交——
检验 $UU^{\mathsf T}$，或在
列的乘积中忘记取共轭，这是辨别错误矩阵
的两种经典办法。
\end{remark}

\begin{example}[等距变换恰好就是酉变换]\label{ex:b2:hermitian:isometry}
保范性看起来比酉性弱，但在 $\C$ 上并非如此：
若 $\norm{u(x)} = \norm x$ 对所有 $x$ 成立，则 $u^*u =
\mathrm{id}$。事实上 $v = u^*u - \mathrm{id}$ 为埃尔米特，
且满足 $\langle x, v(x)\rangle = \norm{u(x)}^2 - \norm x^2 =
0$，对每个 $x$ 成立；由复极化（
上文陷阱 (ii)），“对角线”恒为零的映射必为零：
$v = 0$。具体地，极化过程为
\[
0 = \langle x + y, v(x+y)\rangle
= \langle x, v(y)\rangle + \langle y, v(x)\rangle,
\qquad
0 = \langle x + \iu y, v(x + \iu y)\rangle
= \iu\langle x, v(y)\rangle - \iu\langle y, v(x)\rangle ,
\]
两式合起来即得 $\langle x, v(y)\rangle = 0$，
对所有 $x, y$ 成立。最后一点：这正是“酉”能
只靠测量长度便可检验——这种刚性
正是傅里叶一章将要利用的，那里保持能量
$\norm f_2$（帕塞瓦尔）等价于保持系数的所有
内积。
\end{example}

\begin{remark}[本卷中的几种视角]
这里建立的埃尔米特理论几乎
马上就被用上。傅里叶一章\emph{正是}
无穷维中的埃尔米特几何：指数函数 $(e_n)$
对内积 $\langle f, g\rangle =
\frac{1}{2\pi}\int\conj fg$ 构成标准正交族，贝塞尔不等式
即本章的投影估计，
\cref{thm:b2:hermitian:cs}，而帕塞瓦尔等式
是其极限形式。有限傅里叶变换
（\cref{exo:b2:hermitian:10}）每当卷积
需要对角化时便会重新出现。本章的周末问题——
柯朗–费希尔、外尔、交错定理——提供了
\omterm{def:b2:reduction:eigen}{特征值}的稳定性，这正是微分方程一章
所援引的，断言系统的微小扰动
只会使其频率微小变动。回看前面，这里的一切
都是\omterm{def:b2:quadratic:def}{二次型}一章在复数版本下的镜像：请把
两套对照表并列看待（$A^{\mathsf T} \leftrightarrow
A^\dagger$，正交 $\leftrightarrow$ 酉，\omterm{pb:b2:hermitian:1}{瑞利商}
在两种情形下都是实数）。
\end{remark}

\begin{remark}[正规自同态]
在 $\C$ 上，确切的陈述是：$u$ \emph{酉\omterm{def:b2:reduction:diag}{可对角化}
当且仅当它是正规的}（$u^*u = uu^*$）——
一举涵盖埃尔米特、酉与反埃尔米特映射。这一结论的
证明是对上文论证的一个漂亮加强
（\cref{exo:b2:hermitian:8}）。与之相反，在 $\R$ 上，正规性
只能得到分块对角化（旋转块）：复几何
确实更简单。
\end{remark}

\begin{remark}[应用背景]
埃尔米特谱理论正是量子力学的数学，
可观测量由埃尔米特算子刻画，
（实谱即可测量值），时间演化由酉算子描述
（保范即概率守恒）。
在本书中，傅里叶一章依赖于
指数函数的标准正交性——这是一个\omterm{def:b2:hermitian:def}{埃尔米特内积}的陈述——而
循环矩阵的对角化
（\cref{exo:b2:hermitian:10}）即有限傅里叶变换。
周末问题发展了\omterm{def:b2:reduction:eigen}{特征值}的变分法，
包括柯朗–费希尔、外尔与交错定理，它们是
数值分析与数学物理的日常工具；
本科三年级分册将其推广到希尔伯特空间上的
紧自伴算子。
\end{remark}

\section{练习}

\begin{exercise}[$\star$]\label{exo:b2:hermitian:1}
在 $\C^2$ 中，求 $\langle x, y\rangle$、$\norm x$、$\norm y$，其中
$x = (1, \iu)$，$y = (\iu, 1)$。它们是否正交？给出一组
包含 $\frac{x}{\norm x}$ 的标准正交基。
\end{exercise}

\begin{exercise}[$\star$]\label{exo:b2:hermitian:2}
判断下列矩阵中哪些是埃尔米特的、哪些是酉的、哪些是正规的：
\[
\begin{pmatrix} 1 & \iu\\ -\iu & 2 \end{pmatrix},
\qquad
\frac{1}{\sqrt2}\begin{pmatrix} 1 & \iu\\ \iu & 1\end{pmatrix},
\qquad
\begin{pmatrix} 0 & 1\\ 0 & 0 \end{pmatrix}.
\]
\end{exercise}

\begin{exercise}[$\star$]\label{exo:b2:hermitian:3}
证明：每个矩阵 $A \in \mathcal{M}_n(\C)$ 可唯一写成 $A =
H + \iu K$，其中 $H, K$ 为埃尔米特矩阵，\emph{（“实部与虚部
”$H = \frac{A + A^\dagger}{2}$，$K = \frac{A -
A^\dagger}{2\iu}$）}，并且 $A$ 正规当且仅当 $H$ 与 $K$
可交换。
\end{exercise}

\begin{exercise}[$\star\star$]\label{exo:b2:hermitian:4}
在标准正交基下对角化：$A = \begin{pmatrix} 2 & \iu\\
-\iu & 2\end{pmatrix}$，并计算 $A^k$，其中 $k \in \N$。
\end{exercise}

\begin{exercise}[$\star\star$]\label{exo:b2:hermitian:5}
证明 $U(n)$ 是紧的，并证明\omterm{def:b2:reduction:eigen}{特征值}映射是
满射：每个单位模的 $\lambda$ 都是某个酉矩阵的\omterm{def:b2:reduction:eigen}{特征值}。
证明 $\det U \in \mathbb{U}$（单位圆），对 $U \in
U(n)$ 成立。
\end{exercise}

\begin{exercise}[$\star\star$]\label{exo:b2:hermitian:6}
（凯莱变换）设 $H$ 为埃尔米特矩阵。证明 $H + \iu I$
可逆，且 $U = (H - \iu I)(H + \iu I)^{-1}$ 为酉矩阵，
且 $1 \notin \operatorname{Sp}(U)$。\emph{（用谱的方法：在
$H$ 的特征基上，一切都化为数量。）}
\end{exercise}

\begin{exercise}[$\star\star$]\label{exo:b2:hermitian:7}
设 $A$ 为埃尔米特正定矩阵（$\langle x, Ax\rangle > 0$，对
一切 $x \neq 0$），证明 $\operatorname{Sp}(A) \subseteq
\intoo{0}{\infty}$，并证明 $A = B^2$，
其中 $B$ 为埃尔米特正定矩阵，且 $\det A > 0$。
\end{exercise}

\begin{exercise}[$\star\star\star$]\label{exo:b2:hermitian:8}
（正规自同态的谱定理）设 $u$ 为
\omterm{def:b2:hermitian:def}{埃尔米特空间}上的正规自同态。
\begin{enumerate}
  \item 证明 $\norm{u(x)} = \norm{u^*(x)}$ 对一切 $x$ 成立，
        并由此推出 $\ker(u - \lambda) = \ker(u^* - \conj\lambda)$。
  \item 证明 $u$ 的不同\omterm{def:b2:reduction:eigen}{特征值}所对应的\omterm{def:b2:reduction:eigen}{特征子空间}
        两两正交，并证明\omterm{def:b2:reduction:eigen}{特征子空间}的正交补是
        $u$-稳定的。
  \item 用归纳法得出 $u$ 可酉对角化；
        反之亦然。
\end{enumerate}
\end{exercise}

\begin{exercise}[$\star$]\label{exo:b2:hermitian:9}
称满足 $u^* = -u$ 的自同态为\emph{反\omterm{def:b2:hermitian:adjoint}{埃尔米特自同态}}。证明
其\omterm{def:b2:reduction:eigen}{特征值}均为纯虚数，且映射 $u \mapsto \iu
u$ 为从\omterm{def:b2:hermitian:adjoint}{埃尔米特自同态}到反\omterm{def:b2:hermitian:adjoint}{埃尔米特自同态}的双射，
并证明反\omterm{def:b2:hermitian:adjoint}{埃尔米特自同态}
可酉对角化（见\cref{exo:b2:hermitian:8}）。
\end{exercise}

\begin{exercise}[$\star\star$]\label{exo:b2:hermitian:10}
（有限傅里叶变换）设 $S$ 为
$\C^n$ 上的循环移位：$S(x_0, x_1, \dots, x_{n-1}) = (x_{n-1}, x_0, \dots,
x_{n-2})$，且 $\omega = \eu^{2\iu\pi/n}$。
\begin{enumerate}
  \item 证明 $S$ 为酉变换，且向量 $f_k =
        \frac{1}{\sqrt n}\bigl(1, \omega^k, \omega^{2k}, \dots,
        \omega^{(n-1)k}\bigr)$，$0 \leq k < n$，构成
        由\omterm{def:b2:reduction:eigen}{特征向量}组成的标准正交基：$Sf_k = \omega^{-k}
        f_k$。
  \item 由此推出，每个\emph{循环矩阵} $C =
        \sum_{j=0}^{n-1} c_jS^j$ 为正规的，并可由
        同一组基对角化，\omterm{def:b2:reduction:eigen}{特征值}为 $\widehat c(k) = \sum_j
        c_j\,\omega^{-jk}$。
\end{enumerate}
\end{exercise}

\begin{exercise}[$\star\star$]\label{exo:b2:hermitian:11}
设 $P$ 为\omterm{def:b2:hermitian:def}{埃尔米特空间}上的幂等（$P^2 = P$）自同态。
证明 $P$ 是到其像 $\operatorname{im} P$ 上的\emph{正交}投影，
当且仅当 $P^* = P$。给出
到 $\C v$（$\norm v = 1$）上的正交投影的矩阵，以及
到具有标准正交基 $(v_1, \dots, v_k)$ 的子空间上的正交投影矩阵。
\end{exercise}

\begin{exercise}[$\star\star\star$]\label{exo:b2:hermitian:12}
（用插值法求谱投影）设 $A$ 为埃尔米特矩阵，其
\emph{互异}\omterm{def:b2:reduction:eigen}{特征值}为 $\lambda_1, \dots, \lambda_p$，且
\omterm{def:b2:reduction:eigen}{特征子空间}的直和分解 $E = \bigoplus_i E_i$。定义
拉格朗日多项式 $L_i(X) = \prod_{j\neq i}\frac{X -
\lambda_j}{\lambda_i - \lambda_j}$。证明 $P_i = L_i(A)$
是到 $E_i$ 上的正交投影，且 $P_iP_j = 0$，
当 $i \neq j$ 时，$\sum_i P_i = I$，且 $A = \sum_i \lambda_iP_i$
（称为\emph{谱分解}）；求 $f(A)$ 的表达式，
其中 $f$ 为任一多项式，用 $P_i$ 表示。
\end{exercise}

\section{问题：柯朗–费希尔、外尔与特征值的
计算}

\begin{problem}\label{pb:b2:hermitian:1}
埃尔米特矩阵的\omterm{def:b2:reduction:eigen}{特征值}不只是多项式的根，
它们还是\emph{最优化问题}的解。
这一变分观点——\omterm{pb:b2:hermitian:1}{瑞利商}与
\emph{柯朗–费希尔极小极大定理}\index{柯朗–费希尔
极小极大定理}——使\omterm{def:b2:reduction:eigen}{特征值}可比较、稳定，
并且可计算，本题将汇集它的几项经典成果：
\emph{外尔扰动不等式}\index{外尔扰动不等式
}、\emph{柯西交错}、舒尔与范凯
迹不等式、矩阵平方根的单调性、
以及离散拉普拉斯算子的谱。全文中，$A, B,
E$ 均为 $E = \C^n$ 上的埃尔米特矩阵，其\omterm{def:b2:reduction:eigen}{特征值}按
递减次序列出：$\lambda_1(A) \geq \dots \geq \lambda_n(A)$，
且 $R_A(x) = \frac{\langle x, Ax\rangle}{\langle x, x\rangle}$
在 $x \neq 0$ 时，它称为\emph{瑞利商}\index{瑞利商
}。

\medskip
\noindent\textbf{第一部分——\omterm{pb:b2:hermitian:1}{瑞利商}与极小极大。}
取定一组标准正交的特征基 $(e_1, \dots, e_n)$，满足 $Ae_i =
\lambda_ie_i$。
\begin{enumerate}
  \item 证明 $R_A(x)$ 为实数，并证明
        \[
        \lambda_n \leq R_A(x) \leq \lambda_1
        \qquad (x \neq 0),
        \]
        两个界均可取到：$\lambda_1 = \max R_A$，
        $\lambda_n = \min R_A$。
  \item 证明 $R_A$ 的临界点恰为
        $A$ 的\omterm{def:b2:reduction:eigen}{特征向量}\emph{（对 $t \mapsto R_A(x +
        tv)$ 在 $t = 0$ 处展开，$v$ 任取，再将 $v$ 换为
        $\iu v$）}。
  \item 记 $V_k = \operatorname{Vect}(e_1, \dots, e_k)$ 及
        $W_k = \operatorname{Vect}(e_k, \dots, e_n)$。证明
        \[
        \min_{x \in V_k\setminus\{0\}} R_A(x) = \lambda_k
        = \max_{x \in W_k\setminus\{0\}} R_A(x) .
        \]
  \item 证明\emph{柯朗–费希尔定理}：对 $1 \leq
        k \leq n$，
        \[
        \lambda_k = \max_{\dim V = k}\;\min_{x \in
        V\setminus\{0\}} R_A(x)
        = \min_{\dim W = n-k+1}\;\max_{x \in
        W\setminus\{0\}} R_A(x)
        \]
        \emph{（对任一 $V$，其维数为 $k$：$V \cap W_k \neq
        \{0\}$（由格拉斯曼公式），故 $\min_V R_A \leq \lambda_k$；
        第 3 问表明该界可以取到）}。
  \item （单调性）记 $A \leq B$，当 $B - A$
        为半正定时，由第 4 问可得：$A \leq
        B$ 蕴含 $\lambda_k(A) \leq \lambda_k(B)$，对每个
        $k$ 成立。
\end{enumerate}

\medskip
\noindent\textbf{第二部分——外尔不等式。}
\begin{enumerate}[resume]
  \item 证明：若子空间 $V, W \subseteq \C^n$ 满足 $\dim V
        + \dim W > n$，则两者的交集不为零，并推广之：
        $\dim(V_1 \cap V_2 \cap V_3) \geq \dim V_1 + \dim V_2 +
        \dim V_3 - 2n$。
  \item 证明\emph{外尔不等式}：当 $i + j - 1 \leq n$ 时，
        \[
        \lambda_{i+j-1}(A + B) \leq \lambda_i(A) +
        \lambda_j(B)
        \]
        \emph{（取第 3 问中的子空间 $W_i(A)$、$W_j(B)$ 和
        $V_{i+j-1}(A+B)$，并计算维数
        ）}。
  \item 定义 $\vertiii{E}_2 = \max_{\norm x = 1}\norm{Ex}$，
        并证明 $\vertiii E_2 = \max_k\abs{\lambda_k(E)}$ 对埃尔米特矩阵
        $E$ 成立。由此推出\emph{外尔扰动
        定理}：
        \[
        \bigl|\lambda_k(A + E) - \lambda_k(A)\bigr| \leq
        \vertiii{E}_2
        \qquad (1 \leq k \leq n) :
        \]
        每个\omterm{def:b2:reduction:eigen}{特征值}都是矩阵的 $1$-利普希茨
        函数。
  \item （秩一扰动）设 $P$ 为埃尔米特
        半正定、秩为 $1$ 的矩阵。证明
        \[
        \lambda_k(A) \leq \lambda_k(A + P) \leq
        \lambda_{k-1}(A) \qquad (2 \leq k \leq n),
        \]
        连同 $\lambda_1(A) \leq \lambda_1(A+P)$ 一起成立：新的
        \omterm{def:b2:reduction:eigen}{特征值}与旧的\omterm{def:b2:reduction:eigen}{特征值}\emph{交错}。
  \item 数值验证第 8 问：$A = \begin{pmatrix} 2 &
        \iu\\ -\iu & 2\end{pmatrix}$
        （\cref{exo:b2:hermitian:4}：谱为 $\{3, 1\}$），以及
        $E = \begin{pmatrix} 0 & 1\\ 1 & 0\end{pmatrix}$
        （谱为 $\{1, -1\}$）：计算 $A +
        E$ 的谱，并计算不等式两边的值。
\end{enumerate}

\medskip
\noindent\textbf{第三部分——交错与迹
不等式。}
\begin{enumerate}[resume]
  \item （柯西交错定理）设 $B$ 为 $(n-1)
        \times(n-1)$ 阶、取自 $A$ 左上角的主子矩阵。证明
        \[
        \lambda_{k+1}(A) \leq \lambda_k(B) \leq \lambda_k(A)
        \qquad (1 \leq k \leq n-1)
        \]
        \emph{（视 $\C^{n-1} \subseteq \C^n$；在其上，$R_B$ 是
        $R_A$ 的限制；在两个层次上分别应用柯朗–费希尔极小极大
        定理）}。
  \item 迭代：设 $B$ 为阶为 $n - m$ 的主子矩阵，
        $\lambda_{k+m}(A) \leq \lambda_k(B) \leq
        \lambda_k(A)$。
  \item （舒尔）设 $d_1 \geq d_2 \geq \dots \geq d_n$ 为
        $A$ 的对角元（已排序）。对每个 $k$ 证明：
        \[
        \sum_{i=1}^{k} d_i \leq \sum_{i=1}^{k}\lambda_i(A),
        \]
        且当 $k = n$（即迹）时取等号
        \emph{（所选的 $k$ 个对角元构成一个
        $k\times k$ 阶主子矩阵，由
        第 12 问估计其迹）}。
  \item （凯·范）证明
        \[
        \sum_{i=1}^{k}\lambda_i(A)
        = \max\Bigl\{\sum_{i=1}^{k}\langle x_i, Ax_i\rangle
        : (x_1, \dots, x_k) \text{ 标准正交}\Bigr\} .
        \]
  \item 在下列矩阵上验证第 11 问与第 13 问：
        \[
        A = \begin{pmatrix} 2 & 1 & 0\\ 1 & 2 & 1\\
        0 & 1 & 2\end{pmatrix}
        \qquad
        \bigl(\operatorname{Sp} = \{2 - \sqrt2,\ 2,\
        2+\sqrt2\}\bigr)
        \]
        与其左上角的 $2\times2$ 子块（谱为 $\{1,
        3\}$）及其对角元作对照。
\end{enumerate}

\medskip
\noindent\textbf{第四部分——勒夫纳序。} $A \leq B$
仍表示 $B - A$ 半正定；本部分中的矩阵均为
埃尔米特矩阵。
\begin{enumerate}[resume]
  \item 证明：若 $A \leq B$，则 $a_{ii} \leq b_{ii}$ 对所有
        $i$，有 $\operatorname{tr} A \leq \operatorname{tr} B$，
        并且 $C^\dagger AC \leq C^\dagger BC$ 对\emph{任一}
        复矩阵 $C$ 成立。
  \item 证明平方运算\emph{不}保持单调性：取
        \[
        A = \begin{pmatrix} 1 & 0\\ 0 & 0\end{pmatrix},
        \qquad
        B = \begin{pmatrix} 2 & 1\\ 1 & 1\end{pmatrix},
        \]
        验证 $0 \leq A \leq B$，但 $A^2 \not\leq B^2$。
  \item 证明平方根\emph{是}单调的：若 $0 \leq
        A \leq B$，则 $\sqrt A \leq \sqrt B$ \emph{（设
        $\mu$ 为 $\sqrt B - \sqrt A$ 的一个\omterm{def:b2:reduction:eigen}{特征值}，其单位\omterm{def:b2:reduction:eigen}{特征向量}为
        $v$；计算 $\langle v, (B - A)v\rangle =
        \mu\bigl(\langle v, \sqrt B\,v\rangle + \langle v,
        \sqrt A\,v\rangle\bigr)$，并加以讨论）}。
  \item 证明：求逆在正定矩阵上是反单调的：
        若 $0 < A \leq B$，则 $B^{-1} \leq A^{-1}$
        \emph{（以 $A^{-1/2}$ 作\omterm{def:b2:quadratic:def}{合同}变换，化为 $I \leq M
        \Rightarrow M^{-1} \leq I$，它在谱基下为数的
        不等式）}。
  \item 设 $A, B$ 为正定矩阵。证明
        $AB$ 的\omterm{def:b2:reduction:eigen}{特征值}（一般并非埃尔米特矩阵！）均为
        实数且为正，并且
        \[
        \lambda_{\max}(AB) \leq
        \lambda_{\max}(A)\,\lambda_{\max}(B)
        \]
        \emph{（以 $\sqrt A$ 作相似变换：$AB \sim \sqrt
        A\,B\sqrt A$）}。
\end{enumerate}

\medskip
\noindent\textbf{第五部分——离散拉普拉斯算子的计算示例。}
设 $T_n$ 为 $n \times n$ 三对角矩阵，其主对角线上元素为 $2$，
两条相邻对角线上元素均为 $-1$。
\begin{enumerate}[resume]
  \item 令 $\theta_k = \frac{k\pi}{n+1}$，验证
        向量 $v_k = \bigl(\sin(j\theta_k)\bigr)_{1\leq j\leq
        n}$ 满足 $T_nv_k = (2 - 2\cos\theta_k)\,v_k$
        \emph{（积化和差恒等式；检验边界行
        $j = 1, n$）}。由此得出：
        \[
        \operatorname{Sp}(T_n) = \Bigl\{4\sin^2
        \frac{k\pi}{2(n+1)} : 1 \leq k \leq n\Bigr\},
        \]
        均为单\omterm{def:b2:reduction:eigen}{特征值}，且全部为正，因此 $T_n$ 正定。
  \item （一个势）设实对角矩阵 $D =
        \operatorname{diag}(d_1, \dots, d_n)$，夹逼其谱：
        对每个 $k$，
        \[
        \lambda_k(T_n) + \min_i d_i \;\leq\;
        \lambda_k(T_n + D) \;\leq\; \lambda_k(T_n) + \max_i
        d_i .
        \]
  \item 显式检验 $T_3$ 与 $T_2$ 之间的柯西交错定理
        （谱分别为 $\{2 \pm \sqrt2, 2\}$ 与 $\{1,
        3\}$），并解释：$T_2$ 是由 $T_3$
        去掉路径的一个端点后得到的矩阵。
  \item 证明当 $n \to
        \infty$ 时，极端\omterm{def:b2:reduction:eigen}{特征值}满足：
        \[
        \lambda_{\min}(T_n) = 4\sin^2\frac{\pi}{2(n+1)}
        \sim \frac{\pi^2}{(n+1)^2},
        \qquad
        \lambda_{\max}(T_n) \to 4 ,
        \]
        于是条件数 $\kappa_n =
        \lambda_{\max}/\lambda_{\min}$ 的增长阶为
        $\frac{4(n+1)^2}{\pi^2}$：在越来越细的网格上
        离散化二阶导数，
        本质上是病态的。
  \item 综合。每项用一句话作答：(i) 为何
        使\omterm{def:b2:reduction:eigen}{特征值}稳定的是变分刻画，而不是
        \omterm{def:b2:reduction:charpoly}{特征多项式}；
        （问题 8--9）；(ii) 哪些问题只用了
        $\lambda_1 = \max R_A$，哪些需要完整的
        极小极大刻画；(iii) 勒夫纳序为这一图景增添了什么；
        (iv) 这些工具在何处再次出现（问题 24 刚度矩阵的数值
        分析；量子
        微扰理论；以及在本科三年级分册中，
        紧自伴
        算子的极小极大原理）。
\end{enumerate}
\end{problem}
